% Capitulo: Dados como contexto para IA
\chapter{Dados como contexto para IA}
\label{cap:contexto-ia}

O uso mais comum de IA com dados nas empresas não é treinar modelos. É dar contexto: no
momento da pergunta, o assistente recebe as definições, os documentos e os dados de que
precisa para responder. A qualidade da resposta passa a depender de como esse contexto está
organizado, e organizar contexto é, no fundo, modelar e documentar dados para um consumidor
novo. Este capítulo fecha a Parte~\ref{parte:modelagem} tratando desse consumidor: o que ele
precisa, onde ele erra e como preparar os dados para ele.

\section{O modelo só sabe o que recebe}
\label{sec:modelo-so-sabe}

Um modelo de linguagem sabe muito sobre o mundo em geral e nada sobre a sua empresa. Ele não
conhece as tabelas, as métricas, as siglas internas nem os documentos de procedimento. Quando
recebe uma pergunta sobre esses assuntos, responde com o que está no contexto que recebeu
junto com a pergunta. Quando o contexto não basta, ele não costuma dizer que não sabe: completa
a lacuna com o que parece plausível, e responde com a mesma segurança de sempre.

Há duas formas principais de dar contexto. A primeira é buscar, a cada pergunta, os trechos de
documentos mais relevantes e entregá-los ao modelo junto com a pergunta. A técnica se chama
geração aumentada por recuperação, ou RAG, e foi proposta em 2020 como forma de combinar o
conhecimento geral do modelo com uma base de documentos consultada na hora
\parencite{lewis2020}. A segunda é dar ao modelo acesso a dados estruturados (tabelas,
métricas, consultas) junto com as definições de negócio de que ele precisa para usá-los
corretamente. Padrões abertos para conectar assistentes a fontes de dados e ferramentas, como
o Model Context Protocol \parencite{anthropic2024mcp}, tornaram essa segunda forma muito mais
fácil de montar.

Uma tentação comum é resolver tudo colocando mais contexto: a documentação inteira, todas as
tabelas, todo o histórico. Não funciona tão bem quanto parece. Estudos mostraram que modelos de
linguagem usam pior a informação que fica no meio de contextos muito longos do que a que fica
no começo ou no fim \parencite{liu2024}. E contexto grande custa caro, porque cada token
enviado é cobrado (capítulo~\ref{cap:custo-consumo}). O objetivo é dar o contexto certo, não o
maior possível.

\section{Metadados e camada semântica como contexto}
\label{sec:metadados-contexto}

Para responder perguntas sobre dados estruturados, o assistente precisa do que um analista
novo precisaria: a descrição de cada tabela e coluna, a definição de cada métrica, como as
tabelas se relacionam, quais filtros aplicar por padrão e alguns exemplos de perguntas com a
consulta correta. Uma definição de métrica escrita pensando nesse leitor pode ser simples
assim:

\begin{terminal}
metrica: receita_liquida
descricao: Soma do valor dos itens faturados, menos devoluções e descontos.
  Não inclui frete nem impostos. Usa a data de faturamento.
tabela: fato_vendas
grão: uma linha por item vendido
filtro_padrao: status = 'faturado'
dono: Controladoria
perguntas_exemplo:
  - Qual foi a receita líquida de março?
  - Receita líquida por região no último trimestre
\end{terminal}

Nada nesse trecho é novo para quem trabalha com dados. É a documentação que sempre foi
considerada boa prática e que sempre ficava para depois. A diferença é que agora ela tem um
consumidor que a usa a cada pergunta, e a falta dela aparece imediatamente na forma de
respostas erradas. É o mesmo movimento que transformou os testes em requisito
(\veja{sec:tdd}) e o data steward em revisor (\veja{sec:steward-tutor}): a IA tornou barato,
e necessário, o que antes não cabia no tempo de ninguém. E a própria IA pode ajudar a
escrever a primeira versão dessas descrições, desde que alguém que conhece o negócio revise.

Um modelo dimensional bem feito já é meio caminho andado. Grão explícito, nomes claros,
métricas definidas uma vez numa camada semântica (\veja{sec:camada-semantica}): tudo isso
reduz a chance de o assistente somar a coisa errada ou escolher a coluna de data errada. O
exemplo da receita com frete multiplicado (\veja{sec:consulta-errada}) é exatamente o tipo
de erro que um grão bem documentado ajuda a evitar.

\section{Documentos e recuperação}
\label{sec:documentos-recuperacao}

Para responder com base em documentos (procedimentos, contratos, políticas, relatórios), os
documentos precisam ser preparados. Eles são divididos em trechos de tamanho adequado, cada
trecho recebe metadados (de que documento veio, de que data, quem é o dono, quem pode ler) e
é indexado de forma que a busca encontre os trechos pelo significado, e não só pelas
palavras exatas. Essa busca por significado usa representações numéricas dos textos
(\emph{embeddings}), em que textos parecidos ficam próximos.

O raciocínio da arquitetura medalhão serve aqui também (Figura~\ref{fig:medalhao}). O
documento original fica guardado como chegou, como no bronze. A versão limpa e dividida em
trechos, com os metadados, faz o papel da prata. O índice pronto para a consulta do
assistente é o ouro. Separar as camadas permite refazer a divisão ou trocar a forma de
indexar sem voltar a coletar os documentos.

Dois cuidados são decisivos. O primeiro é a permissão: o assistente não pode responder a uma
pessoa com base num documento que ela não poderia ler. As permissões precisam vir da fonte e
ser respeitadas na busca, e não aplicadas depois, na resposta. O segundo é a atualização:
contexto velho gera resposta velha com cara de nova. Se a política mudou e o índice não foi
atualizado, o assistente vai responder com a regra antiga, com toda a segurança. Tudo o que o
capítulo~\ref{cap:frequencia} discute sobre frequência de atualização vale para o contexto da
IA.

\section{Organizar para a IA}
\label{sec:organizar-para-ia}

Juntando o que foi dito, organizar os dados para servirem de contexto à IA se resume a uma
lista curta, e quase toda ela já apareceu neste livro com outros nomes:

\begin{itemize}
  \item um catálogo em que cada fonte tem dono e descrição;
  \item um glossário de negócio, com os termos e siglas internos;
  \item métricas definidas uma vez, com grão, filtros e dono;
  \item master data conciliado, para que \enquote{Agro Silva} e \enquote{AGRO SILVA LTDA}
    sejam reconhecidos como o mesmo cliente (capítulo~\ref{cap:master-data});
  \item exemplos de perguntas com a resposta ou a consulta correta;
  \item permissões herdadas da fonte;
  \item frequência de atualização definida para cada fonte de contexto.
\end{itemize}

E, para saber se tudo isso funciona, um conjunto de avaliação: perguntas com a resposta
certa conhecida, rodadas a cada mudança de dado, de contexto ou de modelo
(\veja{sec:consumidor-modelo}). Sem ele, cada ajuste é feito no escuro, e uma melhoria numa
pergunta pode piorar outras dez sem que ninguém perceba.

\begin{destaque}
Um assistente que responde perguntas sobre os dados da empresa é tão bom quanto a
documentação e a modelagem que estão por trás dele.
\end{destaque}

\rascunho{
  Relato: o que eu precisei organizar para uma IA responder bem perguntas sobre os nossos
  dados (descrições, métricas, exemplos), o que deu errado no começo e o que mais fez
  diferença. Sem nomes.
}

Com isso, a Parte~\ref{parte:modelagem} se fecha. Ela mostrou que o mesmo dado precisa ser
organizado de formas diferentes para usos diferentes, e que os cadastros, a governança e a
documentação são a base comum de todos eles. A próxima parte trata de como construir e operar
tudo isso com a disciplina da engenharia de software.

\begin{leituras}
  \leitura{lewis2020}{O artigo que propôs a geração aumentada por recuperação (RAG),
    base da maior parte dos assistentes que respondem com documentos da empresa.}
  \leitura{liu2024}{O estudo que mostra por que mais contexto nem sempre é melhor: modelos
    usam pior a informação no meio de contextos longos.}
  \leitura{huyen2022}{Para a visão de sistema: como dados, modelo e avaliação se
    encaixam em produção.}
\end{leituras}
